\section{Numerical Model of EIT Cooling}
\label{sec:numerical_eit}

The EIT cooling dynamics were simulated using QuTiP by solving the open-system evolution of an atom coupled to a single quantized motional mode. Two models of increasing complexity were implemented. The first was an ideal three-level $\Lambda$ system, used to validate the EIT cooling mechanism and the numerical implementation against the analytical predictions introduced in Sec.~\ref{sec:EIT_cooling}. The second model retained the complete hyperfine and Zeeman structure of the $^{87}\mathrm{Rb}$ $D_2$ line and was used to assess the effects of off-resonant excitation, spontaneous-emission branching, optical pumping, and repumping under experimentally relevant conditions.

Both descriptions were formulated within the same Lindblad master-equation framework. They differed only in the dimension and structure of the internal atomic Hilbert space. The motion was treated in one spatial dimension throughout this work. The simulated direction was selected by assigning to each laser the projection of its wavevector along the corresponding trap axis.

\subsection{Internal Atomic Basis}
\label{subsec:atomic_basis}

In the reduced model, the internal Hilbert space was spanned by two ground states, $|g_1\rangle$ and $|g_2\rangle$, and one excited state, $|e\rangle$. A weak probe field with Rabi frequency $\Omega_p$ coupled $|g_1\rangle$ to $|e\rangle$, while a stronger coupling field with Rabi frequency $\Omega_c$ coupled $|g_2\rangle$ to $|e\rangle$. At two-photon resonance, this model supports the dark state of Eq.~\eqref{eq:darkstate} and therefore reproduces the ideal carrier suppression discussed in Sec.~\ref{sec:EIT_cooling}. The three-level calculation was used to test the Hamiltonian construction, the motional coupling, the collapse operators, and the time-evolution routines before introducing the complete atomic structure.

The multilevel model included all hyperfine and Zeeman sublevels of the $5S_{1/2}$ and $5P_{3/2}$ manifolds of $^{87}\mathrm{Rb}$. The internal basis was written as $|F,m_F\rangle$ and comprised the ground manifolds
\begin{equation}
    F=1,\qquad F=2,
\end{equation}
containing three and five Zeeman states, respectively, and the excited manifolds
\begin{equation}
    F'=0,1,2,3,
\end{equation}
containing a total of sixteen states. The dimension of the internal Hilbert space was therefore
\begin{equation}
    \dim\!\left(\mathcal H_{\mathrm{at}}\right)=3+5+16=24.
\end{equation}

All dipole-allowed couplings were constructed explicitly. For a transition between a ground state $|F_g,m_g\rangle$ and an excited state $|F_e,m_e\rangle$, driven by a spherical polarization component $q\in\{-1,0,+1\}$, the coupling strength was calculated from the corresponding hyperfine dipole matrix element. It was written in the form
\begin{equation}
    d_{eg}^{(q)}
    =
    \langle F_e,m_e|d_q|F_g,m_g\rangle
    =
    \langle F_e\Vert d\Vert F_g\rangle
    (-1)^{F_e-m_e}
    \begin{pmatrix}
        F_e & 1 & F_g\\
        -m_e & q & m_g
    \end{pmatrix},
    \label{eq:dipole_matrix_multilevel}
\end{equation}
where the reduced hyperfine matrix element contains the appropriate Wigner-$6j$ coefficient and the reduced fine-structure matrix element. This construction accounts both for the Zeeman Clebsch--Gordan factors and for the relative strengths of transitions involving different hyperfine manifolds. Transitions that do not satisfy the electric-dipole selection rules were assigned zero coupling.

To reduce optical pumping into states outside the EIT cycle, the target $\Lambda$ system was built around states near the edge of a Zeeman manifold rather than around $m_F=0$. In particular, configurations involving $m_F=-2$ were considered because the condition $\Delta m_F=0,\pm1$ restricts the number of ground states accessible by spontaneous decay. Two level configurations were implemented. In the first, both EIT ground states belonged to the $F=2$ manifold, while a repump laser returned population lost to $F=1$. In the second, the two EIT ground states belonged to different hyperfine manifolds, $F=1$ and $F=2$, and the repump field recovered population accumulated in unwanted states outside this pair. The latter configuration was subsequently used for the simulation with the axial trap parameters.

% INSERT FIGURE: Level schemes of the two multilevel configurations.
% The figure should identify probe, coupling, and repump transitions,
% their polarizations, and the dominant spontaneous-emission channels.

\subsection{Motional Hilbert Space}
\label{subsec:motional_basis}

The atomic motion along the selected trap axis was approximated by a harmonic oscillator of angular frequency $\nu$. The total Hilbert space was
\begin{equation}
    \mathcal H
    =
    \mathcal H_{\mathrm{at}}
    \otimes
    \mathcal H_{\mathrm{mot}},
    \label{eq:total_hilbert_space}
\end{equation}
where
\begin{equation}
    \mathcal H_{\mathrm{mot}}
    =
    \mathrm{span}\{|0\rangle,\ldots,|N_{\mathrm{vib}}-1\rangle\}.
\end{equation}
The motional Hamiltonian was
\begin{equation}
    H_{\mathrm{mot}}
    =
    \hbar\nu\left(a^\dagger a+\frac{1}{2}\right),
    \label{eq:motional_hamiltonian_num}
\end{equation}
with $a$ and $a^\dagger$ the annihilation and creation operators. The constant zero-point contribution was omitted in the numerical implementation because it does not affect the dynamics.

The truncation $N_{\mathrm{vib}}$ was chosen by balancing numerical cost against convergence of the motional distribution. For each time-dependent calculation, convergence was assessed by verifying that the population of the highest retained Fock states remained negligible and that increasing $N_{\mathrm{vib}}$ did not appreciably change the calculated mean occupation,
\begin{equation}
    \langle n(t)\rangle
    =
    \mathrm{Tr}\!\left[
        \rho(t)
        \left(\mathbb I_{\mathrm{at}}\otimes a^\dagger a\right)
    \right].
    \label{eq:n_expectation_numerical}
\end{equation}

\subsection{Laser Coupling and Lamb--Dicke Expansion}
\label{subsec:laser_coupling}

For a laser field $l$ with wavevector $\mathbf{k}_l$, optical phase $\phi_l$, and polarization components $\epsilon_{l,q}$, the position-dependent rotating-wave interaction was written as
\begin{equation}
    H_l
    =
    \frac{\hbar}{2}
    \sum_{g,e,q}
    \Omega_l\epsilon_{l,q}\,
    \mathcal C_{eg}^{(q)}
    |e\rangle\langle g|\,
    e^{i k_{l,\parallel}x+i\phi_l}
    +
    \mathrm{h.c.},
    \label{eq:general_laser_interaction}
\end{equation}
where $\mathcal C_{eg}^{(q)}$ denotes the normalized transition amplitude obtained from Eq.~\eqref{eq:dipole_matrix_multilevel}, and
\begin{equation}
    k_{l,\parallel}
    =
    \mathbf{k}_l\cdot\hat{\mathbf e}_{\mathrm{mot}}
\end{equation}
is the projection of the laser wavevector along the simulated motional axis.

Using
\begin{equation}
    x=x_0(a+a^\dagger),
    \qquad
    x_0=\sqrt{\frac{\hbar}{2m\nu}},
\end{equation}
the Lamb--Dicke parameter associated with field $l$ was
\begin{equation}
    \eta_l=k_{l,\parallel}x_0.
    \label{eq:eta_individual_beam}
\end{equation}
In the Lamb--Dicke regime, the position-dependent phase was expanded to first order,
\begin{equation}
    e^{ik_{l,\parallel}x}
    \simeq
    1+i\eta_l(a+a^\dagger).
    \label{eq:LD_expansion_numerical}
\end{equation}
This expansion retains the carrier and first motional sidebands while neglecting higher-order processes.

For a Raman process involving absorption from one field and stimulated emission into another, the relevant two-photon momentum transfer is determined by
\begin{equation}
    \Delta\mathbf{k}
    =
    \mathbf{k}_p-\mathbf{k}_c,
\end{equation}
in agreement with the discussion in Sec.~\ref{sec:EIT_cooling}. The corresponding effective Lamb--Dicke parameter along the simulated direction is
\begin{equation}
    \eta_{\mathrm{eff}}
    =
    (\mathbf{k}_p-\mathbf{k}_c)\cdot
    \hat{\mathbf e}_{\mathrm{mot}}\,x_0.
    \label{eq:eta_effective_numerical}
\end{equation}
In the three-level benchmark, the coupling beam was taken to propagate orthogonally to the motional axis, so that $\eta_c=0$, while the probe field carried the motional coupling. In the multilevel model, an independent value of $\eta_l$ could be assigned to each beam according to the experimental geometry.

The first-order expansion is valid when the occupied motional states satisfy
\begin{equation}
    \eta_l\sqrt{2\langle n\rangle+1}\ll1.
    \label{eq:LD_validity_numerical}
\end{equation}
Calculations initialized outside this regime were treated only as qualitative descriptions of the initial transient. In such cases, the linearized model cannot be regarded as quantitatively exact until the motional occupation has decreased sufficiently.

\subsection{Hamiltonian in the Rotating Frame}
\label{subsec:rotating_hamiltonian}

After transformation to a rotating frame and application of the rotating-wave approximation, the Hamiltonian was expressed as
\begin{equation}
    H
    =
    H_{\mathrm{at}}
    +
    H_{\mathrm{mot}}
    +
    H_{\mathrm{Z}}
    +
    \sum_l H_l.
    \label{eq:total_hamiltonian_numerical}
\end{equation}
Here, $H_{\mathrm{at}}$ contains the hyperfine energies and laser detunings in the selected rotating frame, and
\begin{equation}
    H_{\mathrm{Z}}
    =
    \sum_i
    g_{F_i}\mu_B B m_{F_i}
    |i\rangle\langle i|
    \otimes\mathbb I_{\mathrm{mot}}
    \label{eq:zeeman_hamiltonian}
\end{equation}
describes the linear Zeeman shift produced by the quantization field $B$. The field was assumed sufficiently weak that the hyperfine quantum numbers $F$ and $m_F$ remained appropriate labels.

For the three-level system, the implemented Hamiltonian was
\begin{align}
    H_{\Lambda}
    ={}&
    \hbar\Delta_p|g_1\rangle\langle g_1|
    +
    \hbar\Delta_c|g_2\rangle\langle g_2|
    +
    \hbar\nu a^\dagger a
    \nonumber\\
    &+
    \frac{\hbar\Omega_c}{2}
    \left(
        |e\rangle\langle g_2|
        +
        |g_2\rangle\langle e|
    \right)
    \nonumber\\
    &+
    \frac{\hbar\Omega_p}{2}
    \left[
        |e\rangle\langle g_1|
        \left(1+i\eta_p(a+a^\dagger)\right)
        +
        \mathrm{h.c.}
    \right].
    \label{eq:three_level_numerical_hamiltonian}
\end{align}

In the multilevel model, Eq.~\eqref{eq:general_laser_interaction} was constructed separately for the probe, coupling, and repump fields. Each laser was assigned its own Rabi frequency, detuning, polarization, and projected wavevector. The repump field was introduced to return population from ground states outside the intended $\Lambda$ subsystem. Its motional factor was set to the identity in the simulations reported here,
\begin{equation}
    e^{ik_{\mathrm{rep},\parallel}x}\rightarrow
    \mathbb I_{\mathrm{mot}},
    \label{eq:repump_no_motion}
\end{equation}
so that repump-induced motional sidebands were neglected. This approximation reduces the numerical complexity and isolates the repumper's internal-state recycling function. It also neglects recoil heating associated with repump absorption and must therefore be regarded as an optimistic approximation when repump scattering is appreciable.

\subsection{Spontaneous Emission and Master Equation}
\label{subsec:master_equation_numerical}

The density operator evolved according to
\begin{equation}
    \frac{d\rho}{dt}
    =
    -\frac{i}{\hbar}[H,\rho]
    +
    \sum_k
    \left(
        L_k\rho L_k^\dagger
        -
        \frac{1}{2}
        L_k^\dagger L_k\rho
        -
        \frac{1}{2}
        \rho L_k^\dagger L_k
    \right).
    \label{eq:lindblad_numerical}
\end{equation}
For the three-level system, symmetric branching into the two ground states was assumed:
\begin{equation}
    L_1
    =
    \sqrt{\frac{\gamma}{2}}\,
    |g_1\rangle\langle e|
    \otimes\mathbb I_{\mathrm{mot}},
    \qquad
    L_2
    =
    \sqrt{\frac{\gamma}{2}}\,
    |g_2\rangle\langle e|
    \otimes\mathbb I_{\mathrm{mot}}.
    \label{eq:three_level_collapse}
\end{equation}

For the multilevel system, the collapse operator for a decay channel $e\rightarrow g$ was written as
\begin{equation}
    L_{e\rightarrow g}
    =
    \sqrt{\gamma_e b_{eg}}\,
    |g\rangle\langle e|
    \otimes\mathbb I_{\mathrm{mot}},
    \label{eq:multilevel_collapse}
\end{equation}
where $\gamma_e$ is the total decay rate of the excited state and $b_{eg}$ is the normalized branching fraction,
\begin{equation}
    b_{eg}
    =
    \frac{\displaystyle
        \sum_q\left|d_{eg}^{(q)}\right|^2}
    {\displaystyle
        \sum_{g',q}\left|d_{eg'}^{(q)}\right|^2},
    \qquad
    \sum_g b_{eg}=1.
    \label{eq:branching_fraction}
\end{equation}
This normalization ensures that
\begin{equation}
    \sum_g
    L_{e\rightarrow g}^\dagger
    L_{e\rightarrow g}
    =
    \gamma_e|e\rangle\langle e|.
\end{equation}

In the simulations described here, the spontaneous-emission operators acted trivially on the motional subspace. Consequently, the recoil associated with the emitted photon and the corresponding diffusion in phonon number were not explicitly included. The model therefore accounts for motional changes driven by laser absorption and stimulated coupling, but not for recoil heating during spontaneous emission. Within the Lamb--Dicke regime this correction enters at order $\eta^2$, but it can affect the final occupation and should be included when a quantitative comparison with the theoretical limit of Eq.~\eqref{eq:n_min} or with experimental temperatures is required. A recoil-resolved extension can be obtained by replacing the identity in Eq.~\eqref{eq:multilevel_collapse} with an angularly dependent operator $e^{-i\mathbf{k}_{\mathrm{em}}\cdot\mathbf{x}}$ and averaging over the dipole emission pattern.

\subsection{Steady-State Spectroscopy}
\label{subsec:steady_state_spectroscopy}

The EIT spectrum was calculated by scanning the probe detuning while holding the remaining laser parameters fixed. At each detuning, the stationary density operator was obtained from
\begin{equation}
    \mathcal L_{\mathrm{tot}}(\rho_{\mathrm{ss}})=0,
    \qquad
    \mathrm{Tr}(\rho_{\mathrm{ss}})=1,
    \label{eq:steady_state_condition}
\end{equation}
where $\mathcal L_{\mathrm{tot}}$ denotes the full Liouvillian superoperator.

The principal spectral observable was the total excited-state population,
\begin{equation}
    P_e(\Delta_p)
    =
    \mathrm{Tr}
    \left[
        \rho_{\mathrm{ss}}(\Delta_p)
        \sum_{e}
        |e\rangle\langle e|
    \right].
    \label{eq:excited_population_spectrum}
\end{equation}
In the weak-probe regime, this quantity is proportional to the total photon-scattering rate and was therefore used as an absorption proxy. It should not, however, be identified with the imaginary part of the probe susceptibility when the probe is strong or when several driven transitions contribute simultaneously.

For the multilevel calculation, additional projectors were evaluated to distinguish excitation and population leakage:
\begin{align}
    P_{F'=2}
    &=
    \sum_{e\in F'=2}
    |e\rangle\langle e|,
    \\
    P_{F'=3}
    &=
    \sum_{e\in F'=3}
    |e\rangle\langle e|,
    \\
    P_{\mathrm{out}}
    &=
    \sum_{g\notin\mathcal G_{\Lambda}}
    |g\rangle\langle g|,
    \label{eq:diagnostic_projectors}
\end{align}
where $\mathcal G_{\Lambda}$ is the pair of ground states defining the target $\Lambda$ system. These observables were used to identify the dark resonance, quantify off-resonant excitation, and determine whether population accumulated in ground states not directly addressed by the EIT beams.

The carrier and first motional sidebands were located relative to the two-photon resonance using the detuning convention of the numerical Hamiltonian. Their assignment was checked directly from the matrix elements connecting $|g,n\rangle$ to $|e,n\rangle$ and $|e,n\pm1\rangle$, rather than inferred only from the sign of the plotted detuning. This step was necessary because the labels ``red'' and ``blue'' sideband reverse under a change in detuning convention.

\subsection{Time-Dependent Cooling Dynamics}
\label{subsec:time_dynamics_numerical}

Cooling was quantified by propagating an initial atom--motion state and evaluating Eq.~\eqref{eq:n_expectation_numerical}. The initial internal state was selected within the active ground-state manifold. The motional state was specified either as a Fock state or as a thermal mixture,
\begin{equation}
    \rho_{\mathrm{th}}
    =
    \sum_{n=0}^{N_{\mathrm{vib}}-1}
    p_n|n\rangle\langle n|,
    \qquad
    p_n
    =
    \frac{\bar n_0^n}{(1+\bar n_0)^{n+1}},
    \label{eq:thermal_initial_state}
\end{equation}
with the probabilities renormalized after truncation when necessary.

The three-level model was propagated using QuTiP's deterministic master-equation solver \texttt{mesolve}. This method evolves the density operator and gives noise-free ensemble expectation values. For the full multilevel model, the Monte Carlo wavefunction solver \texttt{mcsolve} was used. It propagates stochastic pure-state trajectories interrupted by quantum jumps and estimates an observable as
\begin{equation}
    \langle O(t)\rangle
    \simeq
    \frac{1}{N_{\mathrm{traj}}}
    \sum_{j=1}^{N_{\mathrm{traj}}}
    \langle\psi_j(t)|O|\psi_j(t)\rangle.
    \label{eq:mc_average}
\end{equation}
The statistical uncertainty decreases approximately as $N_{\mathrm{traj}}^{-1/2}$ for independent trajectories.

For an internal dimension $N_{\mathrm{at}}$ and a motional cutoff $N_{\mathrm{vib}}$, the state-vector dimension is
\begin{equation}
    N=N_{\mathrm{at}}N_{\mathrm{vib}}.
\end{equation}
A density matrix contains $N^2$ complex entries, whereas a Monte Carlo trajectory stores a state vector with $N$ entries. The trajectory method therefore substantially reduced memory use for the 24-level calculation. The computational time of either solver also depends on the sparsity of the Hamiltonian, the number of collapse channels, the stiffness of the differential equations, and the number of trajectories; it is therefore not described by memory scaling alone.

For Monte Carlo calculations, full state storage was disabled and only the required expectation values were retained. Independent trajectories were distributed across the available processor cores. Numerical convergence was assessed by increasing the trajectory number and comparing the resulting cooling curves within their statistical fluctuations. The output time grid was chosen independently of the adaptive internal integration steps used by the solver.

\subsection{Units and Conversion to Physical Time}
\label{subsec:numerical_units}

The numerical calculations were performed in units of the natural angular linewidth,
\begin{equation}
    \gamma=1,
\end{equation}
and $\hbar=1$. All detunings, Rabi frequencies, hyperfine splittings, Zeeman shifts, and trap frequencies supplied to the Hamiltonian were therefore divided by $\gamma$. Time was expressed in units of $\gamma^{-1}$.

For the $^{87}\mathrm{Rb}$ $D_2$ transition, the linewidth was taken as
\begin{equation}
    \frac{\gamma}{2\pi}=6.067~\mathrm{MHz},
\end{equation}
so that one simulation time unit corresponds to
\begin{equation}
    \gamma^{-1}
    =
    \frac{1}{2\pi\times6.067~\mathrm{MHz}}
    \simeq
    26.2~\mathrm{ns}.
    \label{eq:time_conversion}
\end{equation}
Accordingly, a dimensionless interval $t_{\mathrm{sim}}$ was converted to physical time through
\begin{equation}
    t_{\mathrm{phys}}
    =
    \frac{t_{\mathrm{sim}}}{\gamma}.
    \label{eq:physical_time_conversion}
\end{equation}
The same angular-frequency convention was used for $\nu$, $\Omega_l$, $\Delta_l$, and Zeeman shifts to avoid factors of $2\pi$ within the Hamiltonian.

\subsection{Validation of the Three-Level Model}
\label{subsec:three_level_validation}

The reduced model was validated in two stages. First, the motional degree of freedom was omitted and the probe detuning was scanned to verify the formation of the EIT transparency minimum at two-photon resonance and the narrow dressed-state resonance displaced by the coupling-induced AC Stark shift of Eq.~\eqref{eq:acstarkshift}. The numerically obtained peak position was compared with
\begin{equation}
    \delta
    =
    \frac{\sqrt{\Delta_c^2+\Omega_c^2}-|\Delta_c|}{2}.
\end{equation}

Second, the motional basis was restored and the coupling intensity was selected to satisfy the EIT cooling condition. For the sign convention and blue-detuned configuration considered in Sec.~\ref{sec:EIT_cooling}, the exact relation is given by Eq.~\eqref{eq:exact_cooling_condition},
\begin{equation}
    \Omega_c^2=4\nu(\nu+\Delta_c),
    \label{eq:exact_condition_numerical}
\end{equation}
which reduces to
\begin{equation}
    \Omega_c^2\simeq4\Delta_c\nu
    \label{eq:approx_condition_numerical}
\end{equation}
when $\nu\ll\Delta_c$. Under this condition, the narrow absorption resonance is aligned with the phonon-removing sideband, while carrier excitation is suppressed by the dark resonance.

The simulated late-time occupation and cooling rate were compared with the analytical rate-equation results of Eqs.~\eqref{eq:mean_phonon_evolution}--\eqref{eq:n_min}. Agreement was expected only in the weak-probe, unsaturated, Lamb--Dicke, and sufficiently large motional-basis limits. The omission of spontaneous-emission recoil from the numerical collapse operators was also taken into account when interpreting deviations from the analytical steady-state limit.

\subsection{Optimization of the Multilevel Configuration}
\label{subsec:multilevel_optimization}

The multilevel laser parameters were optimized using steady-state spectra before performing the more expensive time-dependent simulations. For each parameter set, the probe detuning was scanned and the following quantities were recorded: the total excited-state population, the populations of the target and off-resonant excited manifolds, the population outside the active $\Lambda$ subsystem, the detuning of the transparency minimum, and the detuning of the nearest absorption maximum.

The scans were performed sequentially:

\begin{enumerate}
    \item The repump Rabi frequency was varied to determine the minimum intensity required to prevent long-lived population accumulation outside the EIT cycle. Excessive repump power was avoided because off-resonant coupling can shift and broaden the EIT spectrum.

    \item The probe Rabi frequency was varied within the weak-probe regime. Increasing $\Omega_p$ increases the scattering signal and may increase the cooling rate, but it also saturates and power-broadens the dark resonance, increasing carrier scattering.

    \item The coupling detuning and coupling Rabi frequency were adjusted to align the narrow absorption maximum with the phonon-removing sideband after accounting for the light shifts introduced by the complete multilevel structure and by the repump field.
\end{enumerate}

The optimization criterion was not the height of the absorption peak alone. Parameter sets were selected to provide low carrier excitation, enhanced scattering at the cooling sideband, reduced scattering at the heating sideband, and limited population leakage. The selected parameters were then tested through a time-dependent \texttt{mcsolve} calculation of $\langle n(t)\rangle$. Because the scans were sequential rather than multidimensional, the resulting set should be interpreted as an operational optimum within the explored ranges, rather than as a unique global optimum.

\subsection{Simulation with Axial Experimental Parameters}
\label{subsec:axial_simulation}

The configuration in which the two EIT ground states belong to different hyperfine manifolds was subsequently evaluated using the axial trap parameters of the experimental system. The axial trap frequency was inserted directly into the motional Hamiltonian, while the Lamb--Dicke parameters of the probe and coupling beams were calculated from their wavevector projections according to Eq.~\eqref{eq:eta_individual_beam}. The magnetic field, laser polarizations, hyperfine splittings, detunings, and Rabi frequencies were retained explicitly in the 24-level Hamiltonian.

The steady-state probe scan was repeated for the axial parameters because changing $\nu$ changes the position of the target motional sideband and therefore the required dressed-state shift. The probe, coupling, and repump parameters were adjusted until the absorption maximum was aligned with the axial phonon-removing sideband while the carrier remained close to the transparency minimum. The resulting parameter set was then used in the Monte Carlo time evolution.

For this calculation, convergence was checked with respect to the motional cutoff, the number of trajectories, and the output time interval. The cooling time was extracted by fitting the ensemble-averaged occupation to
\begin{equation}
    \langle n(t)\rangle
    =
    n_{\infty}
    +
    \left(n_0-n_{\infty}\right)e^{-Wt},
    \label{eq:cooling_fit}
\end{equation}
over the interval in which a single-exponential description was adequate. Here, $W$ is the fitted cooling rate and $n_{\infty}$ is the asymptotic mean occupation. When the curve exhibited multiple timescales due to optical pumping among internal states, Eq.~\eqref{eq:cooling_fit} was used only as an effective characterization and the fitted interval was reported together with the result.

\subsection{Adaptation to In-Fiber Cooling}
\label{subsec:in_fiber_numerical_extension}

To extend the model to cooling inside the hollow-core photonic crystal fiber, the spatially dependent differential AC Stark shift derived in Sec.~\ref{sec:EIT_cooling} must be included in the internal Hamiltonian. For a transverse atomic position $\mathbf r$, the trap-induced shift of each internal state is
\begin{equation}
    H_{\mathrm{trap,AC}}(\mathbf r)
    =
    \sum_i
    U_i(\mathbf r)|i\rangle\langle i|,
    \label{eq:trap_ac_hamiltonian}
\end{equation}
with $U_i(\mathbf r)$ determined by the dynamic polarizability at the trapping wavelength. The local optical detuning then becomes
\begin{equation}
    \Delta_l(\mathbf r)
    =
    \Delta_l^{(0)}
    -
    \frac{U_e(\mathbf r)-U_g(\mathbf r)}{\hbar}.
    \label{eq:position_dependent_detuning}
\end{equation}

A direct three-dimensional treatment would require simultaneous quantization of the axial and two transverse motional modes and would substantially increase the Hilbert-space dimension. In the present one-dimensional framework, the leading effect of the transverse motion can instead be incorporated by averaging the local steady-state or sideband rates over the transverse spatial distribution:
\begin{equation}
    \overline{A}_{\pm}
    =
    \int d^2r\,
    p_T(\mathbf r)\,
    A_{\pm}\!\left[\Delta_l(\mathbf r),\Omega_l(\mathbf r),\nu(\mathbf r)\right],
    \label{eq:spatially_averaged_rates}
\end{equation}
where $p_T(\mathbf r)$ is the thermal transverse probability density. In the harmonic approximation,
\begin{equation}
    p_T(r)
    =
    \frac{m\nu_r^2}{2\pi k_BT}
    \exp\!\left(
        -\frac{m\nu_r^2r^2}{2k_BT}
    \right).
    \label{eq:transverse_thermal_distribution}
\end{equation}
This procedure incorporates the inhomogeneous broadening estimated in Eq.~\eqref{eq:inhom_broad} without explicitly adding two further harmonic-oscillator Hilbert spaces.

For each sampled position, the local probe and coupling detunings, Rabi frequencies, and trap frequencies must be recalculated from the fiber-mode intensity. The on-axis laser parameters can then be chosen to satisfy the modified resonance condition discussed in Sec.~\ref{sec:EIT_cooling}, while Eq.~\eqref{eq:spatially_averaged_rates} quantifies the degradation caused by atoms sampling different local light shifts. This extension provides the connection between the free-space multilevel simulation and the in-fiber cooling geometry.

\subsection{Scope and Approximations}
\label{subsec:numerical_approximations}

The numerical model includes the full internal hyperfine and Zeeman structure, coherent coupling by the applied laser fields, first-order motional sidebands for the EIT beams, and all dipole-allowed spontaneous-emission branches. The main approximations are the restriction to one motional dimension, harmonic confinement, the first-order Lamb--Dicke expansion, omission of spontaneous-emission recoil, neglect of repump motional coupling, and sequential rather than global parameter optimization. The idealized three-level results are therefore used primarily as a validation of the EIT mechanism, while the 24-level calculations quantify the internal-state effects relevant to the experimental implementation. Quantitative predictions for the final in-fiber temperature require the spatial averaging of Eq.~\eqref{eq:spatially_averaged_rates}, together with recoil-resolved collapse operators and experimentally calibrated position-dependent light shifts.
